%==========================================================
%  CHAPTER 2: Inflation, Monetary Policy, and the Cost of Living
%  ECON 204 — Contemporary Economic Issues: A Canadian Perspective
%==========================================================

\chapter{Inflation, Monetary Policy, and the Cost of Living}
\label{ch:inflation}

%----------------------------------------------------------
%  OPENING VIGNETTE
%----------------------------------------------------------

\noindent In January 2021, Fatima filled her grocery cart in Saskatoon and
paid \$183. By January 2022, the same cart cost \$206. By June 2022, it
cost \$228. In eighteen months, her weekly grocery bill had risen by
\$45 --- a 25\% increase on food alone. Gasoline at the pump had gone from
\$1.04 per litre to \$1.85. Her landlord had issued a rent increase notice
citing rising property taxes and maintenance costs. Her variable-rate
mortgage payment had climbed \$390 per month since March 2022, tracking the
Bank of Canada's overnight rate almost in lockstep.

Fatima had not changed her lifestyle. She had not moved to a more expensive
neighbourhood or started eating at restaurants. She had the same job, the
same apartment, and the same spending habits. But the purchasing power of
her income was shrinking, month by month, in a way she could not control
and could barely predict.

Statistics Canada would later record that Canada's Consumer Price Index rose
8.1\% in June 2022 --- the highest reading in 40 years. That number meant
something very specific: an average Canadian household needed \$108.10 in
June 2022 to buy what \$100 bought in June 2021. For Fatima, whose budget
had less slack than the average, the real erosion was larger.

What caused inflation this severe? Why did it emerge so suddenly after three
decades of low, stable prices? How does the Bank of Canada fight it, and
at what cost? And who bears the heaviest burden when prices rise faster
than wages? These are the questions this chapter answers.

%----------------------------------------------------------
%  LEARNING OBJECTIVES
%----------------------------------------------------------

\begin{learningobjectives}
  \item Define inflation, calculate the CPI inflation rate from price
        data, and identify the four main biases in CPI measurement.
  \item Distinguish headline CPI from core inflation measures and explain
        why the Bank of Canada monitors both.
  \item Apply the Fisher equation to calculate real interest rates from
        nominal rates and inflation, and derive real wage growth from
        nominal wage and CPI data.
  \item Use the aggregate demand--aggregate supply framework to explain
        demand-pull and cost-push inflation, and apply both to Canada's
        2021--2023 episode.
  \item Describe the short-run and long-run Phillips Curve, interpret
        Canadian data on the inflation--unemployment relationship, and
        define the sacrifice ratio.
  \item Explain how the Bank of Canada's inflation-targeting framework
        works, trace the channels of monetary policy transmission, and
        evaluate the 2022--2024 rate-hike cycle.
  \item Analyse the distributional consequences of inflation across
        income groups, between debtors and creditors, and between
        workers with indexed and non-indexed compensation.
\end{learningobjectives}

%==========================================================
\section{Why Inflation Matters}
%==========================================================

\term{Inflation} is a sustained increase in the general price level of goods
and services in an economy over time. Three words in that definition deserve
emphasis. \textbf{Sustained}: a single price spike --- gasoline surging one
month and falling the next --- is not inflation. \textbf{General}: inflation
refers to the average price level, not any specific good. A single price
rising while others fall is a \term{relative price} change, not inflation.
\textbf{Over time}: inflation is a dynamic phenomenon measured as a rate
of change, typically expressed as a percentage per year.

The opposite of inflation --- a sustained \textit{decrease} in the general
price level --- is \term{deflation}. Deflation may sound attractive
(prices falling means money buys more), but it is one of the most dangerous
macroeconomic conditions an economy can experience. When prices are expected
to fall, households postpone purchases (why buy today what will be cheaper
tomorrow?), firms delay investment, profits shrink, layoffs rise, and the
economy can fall into a self-reinforcing downward spiral. Japan's ``Lost
Decade'' of the 1990s is the modern cautionary example.

Low, stable, \textit{positive} inflation --- in Canada's case, the
2\% target --- is the policy goal. Why 2\% rather than 0\%? Three reasons:

\begin{enumerate}
  \item \textbf{Measurement bias:} The CPI tends to overstate true inflation
        by approximately 0.3--0.6 percentage points because it cannot
        fully capture quality improvements, new goods, and consumer
        substitution toward cheaper alternatives. A 2\% measured target
        therefore implies roughly zero true inflation.
  \item \textbf{Policy buffer:} If the equilibrium real interest rate is
        approximately 2\% and inflation is targeted at 2\%, the Bank of
        Canada can set a nominal rate of approximately 4\% at equilibrium.
        This leaves 400 basis points of room to cut rates before hitting
        the zero lower bound in a recession. A 0\% inflation target would
        leave only 200 basis points of room --- dangerously thin.
  \item \textbf{Wage adjustment:} Inflation allows real wages to fall
        without requiring nominal pay cuts, which workers strongly resist.
        In a world of 2\% inflation, a firm that needs to reduce labour
        costs can simply freeze wages --- real wages fall while nominal
        wages stay constant.
\end{enumerate}

\begin{thinkbox}
If 2\% inflation is the goal, should the Bank of Canada have started
raising rates in late 2021 when inflation was already at 4.7\%, rather than
waiting until March 2022 when it had reached 5.7\%? This is one of the
most debated questions in recent Canadian monetary policy history.
Your answer requires a theory of how quickly monetary policy works
(the transmission lags), an assessment of whether the inflation was
supply-side or demand-side, and a judgement about the risk of
over-tightening into recession. We return to this question in the
Policy Debate at the end of this chapter.
\end{thinkbox}

%==========================================================
\section{Measuring Inflation: The Consumer Price Index}
\label{sec:cpi_measurement}
%==========================================================

%----------------------------------------------------------
\subsection{The CPI Basket Approach}
%----------------------------------------------------------

Canada's primary measure of inflation is the \term{Consumer Price Index (CPI)},
produced monthly by Statistics Canada. The CPI measures the change in the
total cost of a fixed \term{representative basket} of goods and services
purchased by a typical Canadian urban household.

\begin{defbox}{Consumer Price Index (CPI)}
The \textbf{Consumer Price Index} measures the average change over time
in the prices paid by urban consumers for a fixed basket of goods and
services.

\textit{Economic intuition:} The CPI asks: ``How much does it cost today,
relative to a base period, to maintain a fixed standard of living?''
If the basket cost \$1{,}000 in the base period and costs \$1{,}081 today,
the CPI is 108.1, implying 8.1\% cumulative inflation since the base.

\textit{Formal definition:} Statistics Canada periodically determines the
basket weights from the Survey of Household Spending. Prices are collected
from approximately 30{,}000 retail outlets, service providers, and
landlords across Canada each month.

\textit{Canadian note:} Statistics Canada publishes the CPI monthly in
Table 18-10-0004-01. The June 2022 CPI reading was the highest year-over-year
change since September 1983, when Canada was in the final stages of
fighting the 1970s inflation surge.
\end{defbox}

Construction of the CPI proceeds in three steps:

\begin{enumerate}
  \item \textbf{Determine basket weights.} Statistics Canada's Survey of
        Household Spending identifies what Canadian urban households buy
        and in what proportions. These proportions are the
        \term{expenditure weights} $w_i$ for each category $i$.
  \item \textbf{Collect prices.} Each month, approximately 100{,}000 price
        observations are gathered from retailers, landlords, and service
        providers across Canada.
  \item \textbf{Construct the index.} The CPI compares current-period prices
        to base-period prices, weighted by expenditure shares.
\end{enumerate}

The formal CPI formula is a Laspeyres price index:

\begin{equation}
  \text{CPI}_t = \frac{\displaystyle\sum_{i} w_i \cdot P_{i,t}}
                      {\displaystyle\sum_{i} w_i \cdot P_{i,0}}
                 \times 100
  \label{eq:cpi_formula}
\end{equation}

\noindent where $P_{i,t}$ is the price of good $i$ in period $t$,
$P_{i,0}$ is the price of good $i$ in the base period, and $w_i$ is the
expenditure weight of good $i$.

The \term{inflation rate} between two consecutive periods is:

\begin{equation}
  \pi_t = \frac{\text{CPI}_t - \text{CPI}_{t-1}}{\text{CPI}_{t-1}} \times 100
  \label{eq:inflation_rate}
\end{equation}

\noindent\textbf{Worked Example 2.1 --- Calculating the CPI from a Simple Basket}

Suppose a household's basket consists of three goods:

\begin{center}
\renewcommand{\arraystretch}{1.3}
\begin{tabular}{lrrrr}
\toprule
Good & Weight $w_i$ & Base Price $P_{i,0}$ & Current Price $P_{i,t}$
     & Price Relative $P_{i,t}/P_{i,0}$ \\
\midrule
Food       & 0.16 & \$120 & \$131.40 & 1.095 \\
Shelter    & 0.30 & \$900 & \$963.00 & 1.070 \\
Transport  & 0.20 & \$300 & \$342.00 & 1.140 \\
\midrule
\multicolumn{5}{l}{\textit{(Other categories omitted for simplicity)}} \\
\bottomrule
\end{tabular}
\end{center}

\noindent\textit{Step 1:} Calculate the weighted sum of price relatives:
\[
  \text{CPI}_t = (0.16 \times 1.095) + (0.30 \times 1.070)
                + (0.20 \times 1.140) + \ldots
\]
\[
  = 0.1752 + 0.3210 + 0.2280 + \ldots = 0.7242 \text{ (for these three)}
\]

\noindent\textit{Full CPI} (with all categories at assumed price relatives):
Scaling up with remaining categories at an assumed average relative of 1.060:
\[
  \text{CPI}_t \approx 107.5 \implies \pi_t \approx 7.5\%
\]

\noindent\textit{Interpretation:} The basket now costs 7.5\% more than
in the base period. Because shelter (the largest weight) rose only 7.0\%
while transport (a smaller weight) rose 14.0\%, the shelter-weighted
average moderates what would otherwise be a higher measured rate.

\begin{table}[H]
\centering
\caption{Canada CPI Inflation by Component, 2020--2024 (Annual Average, \%)}
\label{tab:cpi_components}
\renewcommand{\arraystretch}{1.30}
\begin{tabular}{lrrrrr}
\toprule
\textbf{CPI Component} & \textbf{2020} & \textbf{2021} & \textbf{2022}
                       & \textbf{2023} & \textbf{2024} \\
\midrule
All-items (headline)  &  0.7 &  3.4 &  6.8 & 3.9 & 2.6 \\
\midrule
Food                  &  2.4 &  3.9 &  9.7 & 5.9 & 1.5 \\
\quad Food purchased from stores & 2.8 & 3.1 & 10.3 & 6.5 & 1.2 \\
\quad Food from restaurants & 1.5 & 5.0 & 8.3 & 5.1 & 2.1 \\
Shelter               &  2.4 &  4.6 &  6.9 & 6.0 & 5.7 \\
\quad Mortgage interest costs & 0.8 & 1.3 & 11.5 & 28.5 & 19.4 \\
\quad Rent             &  3.0 &  1.4 &  4.5 &  8.1 &  8.9 \\
Transportation        & $-$3.6 &  8.3 & 10.6 & 2.4 & 0.8 \\
\quad Gasoline        &$-$14.3 & 32.2 & 31.7 &$-$4.2 &$-$2.1 \\
Energy (total)        & $-$8.5 & 19.5 & 28.5 &$-$8.5 &$-$2.7 \\
Services              &  1.4 &  2.5 &  5.5 & 4.7 & 4.0 \\
Goods (excl.\ food/energy) & 0.3 & 4.3 & 7.2 & 2.7 & 0.5 \\
\bottomrule
\end{tabular}
\source{Statistics Canada, Table 18-10-0004-01, CPI by component. Values
are annual averages; 2024 values are estimates consistent with published data.}
\end{table}

\begin{figure}[H]
\centering
\begin{tikzpicture}
\begin{axis}[
  ybar,
  bar width=0.55cm,
  width=13cm, height=6.5cm,
  symbolic x coords={{Shelter},{Transport},{Food},{Rec./Educ.},
                     {Household},{Health},{Clothing},{Alc./Tob.}},
  xtick=data,
  x tick label style={rotate=35, anchor=east, font=\small},
  ylabel={Weight in CPI Basket (\%)},
  ylabel style={font=\small},
  ymin=0, ymax=34,
  ytick={0,5,10,15,20,25,30},
  yticklabel style={font=\small},
  ymajorgrids=true,
  grid style={dashed, gray!25},
  enlarge x limits=0.08,
  nodes near coords,
  nodes near coords style={font=\footnotesize},
  every axis plot/.append style={fill=unbcgreen!70, draw=unbcgreen!90}
]
\addplot coordinates {
  ({Shelter},    30.0)
  ({Transport},  20.0)
  ({Food},       16.0)
  ({Rec./Educ.}, 11.0)
  ({Household},  10.0)
  ({Health},      5.0)
  ({Clothing},    5.0)
  ({Alc./Tob.},   3.0)
};
\end{axis}
\end{tikzpicture}
\caption{Weights in Canada's CPI Basket, 2022. Shelter and transportation
together account for 50\% of the CPI basket. This structural weighting
means that even moderate increases in rent or mortgage costs --- both
classified under Shelter --- move the headline inflation rate substantially.
The dominance of Shelter also means that the Bank of Canada's rate hikes,
which raised mortgage costs sharply, paradoxically \textit{increased}
one component of the CPI (mortgage interest costs) even as they reduced
the underlying demand driving inflation.}
\label{fig:cpi_basket}
\source{Statistics Canada, CPI Basket Weights, 2022 reference basket.}
\end{figure}

%----------------------------------------------------------
\subsection{Biases in CPI Measurement}
%----------------------------------------------------------

The CPI is a Laspeyres index: it holds the basket fixed at its base-period
composition. This creates systematic upward biases in measured inflation:

\begin{itemize}
  \item \textbf{Substitution bias:} When the price of good A rises relative
        to good B, consumers substitute toward B. The fixed CPI basket
        does not reflect this --- it still assumes the old, pre-increase
        quantities of A are purchased. Result: the CPI overstates the
        true cost-of-living increase.

  \item \textbf{Quality bias:} If a new car costs \$5{,}000 more than last
        year's model but is also significantly better (safer, more
        fuel-efficient, with better technology), not all of the \$5{,}000
        increase is inflation. Some of it is more quality per dollar.
        Statistics Canada adjusts for some quality changes, but imperfectly.

  \item \textbf{New goods bias:} When a completely new product (say,
        smartphones in 2007) enters the market, it is not immediately
        included in the basket. The large initial price declines as new
        products diffuse through the market are missed.

  \item \textbf{Outlet substitution bias:} When consumers shift purchases
        from traditional retailers to lower-priced online sellers, the
        CPI may not fully capture the effective price reduction.
\end{itemize}

Studies by Statistics Canada and the Bank of Canada suggest the CPI
overstates true inflation by approximately 0.3--0.5 percentage points
per year on average. This is one reason the Bank of Canada's 2\% inflation
target is set above zero rather than at zero.

%----------------------------------------------------------
\subsection{Canada's Annual Inflation Record, 2010--2026}
%----------------------------------------------------------

\begin{figure}[H]
\centering
\begin{tikzpicture}
\begin{axis}[
  width=14cm, height=7cm,
  xlabel={Year},
  ylabel={CPI Inflation, Year-over-Year (\%)},
  xlabel style={font=\small},
  ylabel style={font=\small},
  xmin=2009.5, xmax=2026.5,
  ymin=-1, ymax=10,
  xtick={2010,2012,2014,2016,2018,2020,2022,2024,2026},
  xticklabel style={font=\small},
  ytick={-1,0,1,2,3,4,5,6,7,8,9,10},
  yticklabel style={font=\small},
  ymajorgrids=true,
  grid style={dashed, gray!20},
  axis line style={color=gray!60},
  tick style={color=gray!60},
  clip=false
]
% Control range band 1--3%
\fill[unbcgreen!12] (axis cs:2009.5,1) rectangle (axis cs:2026.5,3);
% 2% target line
\addplot[unbcgreen!60, dashed, thick] coordinates {(2009.5,2.0)(2026.5,2.0)}
  node[right, font=\footnotesize, xshift=2pt, color=unbcgreen!80] {2\% target};
% 1% lower bound
\addplot[unbcgreen!40, dotted] coordinates {(2009.5,1.0)(2026.5,1.0)};
% 3% upper bound
\addplot[unbcgreen!40, dotted] coordinates {(2009.5,3.0)(2026.5,3.0)};
% CPI data line
\addplot[policyblue, thick, mark=*, mark size=1.5pt] coordinates {
  (2010, 1.8)(2011, 2.9)(2012, 1.5)(2013, 0.9)(2014, 2.0)
  (2015, 1.1)(2016, 1.4)(2017, 1.6)(2018, 2.3)(2019, 1.9)
  (2020, 0.7)(2021, 3.4)(2022, 6.8)(2023, 3.9)(2024, 2.6)
  (2025, 2.1)(2026, 2.0)
};
% Annotations
\node[font=\footnotesize, align=center, text=policyblue, fill=white,
      inner sep=1pt] at (axis cs:2022,7.5)
  {June 2022 peak: 8.1\%\\(annual avg: 6.8\%)};
\draw[->, policyblue, thin] (axis cs:2022,7.3) -- (axis cs:2022,6.9);
\node[font=\footnotesize, color=gray] at (axis cs:2020,0.2) {COVID};
\node[font=\footnotesize, align=center, color=unbcgreen!70] at (axis cs:2017,0.2)
  {Control range\\(1--3\%)};
\end{axis}
\end{tikzpicture}
\caption{Canada Annual CPI Inflation Rate, 2010--2026. The shaded band
shows the Bank of Canada's 1--3\% control range; the dashed line marks the
2\% midpoint target. For most of 2012--2021, inflation remained within or
near the band, reflecting credible monetary policy. The 2020 drop to 0.7\%
(pandemic demand collapse) and the subsequent surge to a peak of 8.1\%
in June 2022 (annual average 6.8\%) represent the most dramatic deviation
from target in the Bank's 30-year history of inflation targeting. By 2024,
inflation had returned to the upper portion of the target range.}
\label{fig:cpi_annual}
\source{Statistics Canada, Table 18-10-0004-01; Bank of Canada. Annual averages;
2025--2026 are estimates consistent with Bank of Canada projections.}
\end{figure}

\begin{keytakeaway}
The CPI is the standard measure of inflation but contains systematic
upward biases of approximately 0.3--0.5 percentage points per year.
The Bank of Canada targets 2\% measured CPI inflation, which is
roughly consistent with near-zero \textit{true} inflation once biases
are accounted for. Understanding CPI biases is essential for interpreting
economic data and evaluating policy.
\end{keytakeaway}

%==========================================================
\section{Core Inflation: Seeing Through the Noise}
\label{sec:core_inflation}
%==========================================================

Headline CPI measures the price level experienced by the average household.
But for \textit{monetary policy}, headline CPI is a problematic guide. The
reason is that a central bank can only affect demand conditions --- it cannot
directly lower the price of crude oil, ease a shipping container shortage,
or replace a failed harvest. When inflation is driven by temporary supply
shocks (a spike in gasoline prices, a drought-induced food price jump), a
central bank that responds aggressively risks \term{policy error}: tightening
into a supply-driven recession.

\term{Core inflation} measures remove the most volatile components from
the headline CPI to reveal the underlying, persistent inflation trend that
monetary policy can actually address. The Bank of Canada monitors three
official core measures:

\begin{defbox}{Bank of Canada Core Inflation Measures}
\textbf{CPI-trim:} Removes the tails of the monthly price change
distribution --- the highest and lowest 20\% of components by price
change magnitude --- and computes the weighted average of the remaining
60\%.

\textbf{CPI-median:} The price change at the 50th percentile of the
weighted distribution of component price changes. Exactly half of the
CPI basket (by weight) is changing faster than this rate, and half
slower.

\textbf{CPI-common:} Extracts the common variation in price changes
across CPI components using statistical factor analysis. It captures
inflation that is economy-wide rather than sector-specific.

\textit{Why three measures?} Each measure has different statistical
properties. By monitoring all three and looking for consensus, the Bank
reduces the risk of any single measure sending a misleading signal.
When all three core measures are elevated, the Bank is confident that
inflation is broad-based, persistent, and amenable to a policy response.
\end{defbox}

\begin{figure}[H]
\centering
\begin{tikzpicture}
\begin{axis}[
  width=14cm, height=7cm,
  xlabel={Year},
  ylabel={Inflation Rate (\%, Year-over-Year)},
  xlabel style={font=\small},
  ylabel style={font=\small},
  xmin=2017.8, xmax=2026.5,
  ymin=0, ymax=9,
  xtick={2018,2019,2020,2021,2022,2023,2024,2025,2026},
  xticklabel style={font=\small},
  ytick={0,1,2,3,4,5,6,7,8,9},
  yticklabel style={font=\small},
  ymajorgrids=true,
  grid style={dashed, gray!20},
  legend pos=north west,
  legend style={font=\small, fill=white, fill opacity=0.9},
  clip=false
]
% Control range band
\fill[unbcgreen!10] (axis cs:2017.8,1) rectangle (axis cs:2026.5,3);
\addplot[unbcgreen!50, dashed] coordinates {(2017.8,2)(2026.5,2)};
% Headline CPI
\addplot[gray!70, thick, mark=none] coordinates {
  (2018,2.3)(2019,1.9)(2020,0.7)(2021,3.4)(2022,6.8)(2023,3.9)(2024,2.6)(2025,2.1)(2026,2.0)
};
\addlegendentry{Headline CPI}
% CPI-trim (smoother, peaks lower than headline)
\addplot[policyblue, thick, mark=none] coordinates {
  (2018,2.1)(2019,2.1)(2020,1.5)(2021,2.9)(2022,5.5)(2023,3.7)(2024,2.7)(2025,2.2)(2026,2.0)
};
\addlegendentry{CPI-trim}
% CPI-median
\addplot[casered, thick, mark=none, dashed] coordinates {
  (2018,2.0)(2019,2.2)(2020,1.9)(2021,2.7)(2022,5.0)(2023,3.5)(2024,2.6)(2025,2.1)(2026,2.0)
};
\addlegendentry{CPI-median}
% CPI-common
\addplot[dataorange, thick, mark=none, dotted] coordinates {
  (2018,1.9)(2019,1.9)(2020,1.4)(2021,2.2)(2022,4.6)(2023,3.1)(2024,2.5)(2025,2.1)(2026,2.0)
};
\addlegendentry{CPI-common}
% Annotations
\node[font=\footnotesize, text=gray!80] at (axis cs:2022.3,7.2)
  {Headline 8.1\% peak};
\node[font=\footnotesize, text=policyblue, align=center] at (axis cs:2022.8,6.0)
  {Core peaks:\\5--6\%};
\end{axis}
\end{tikzpicture}
\caption{Headline CPI vs.\ Bank of Canada Core Inflation Measures, 2018--2026.
The core measures (CPI-trim, CPI-median, CPI-common) peaked approximately
2--3 percentage points below the 8.1\% headline peak, reflecting the
contribution of volatile energy and some food prices to the headline figure.
Crucially, all three core measures rose substantially above the 3\% upper
bound of the control range in 2022 --- confirming that the inflation surge
was broad-based, not merely a transitory spike in a few volatile categories.
This broad-based elevation justified the Bank of Canada's aggressive tightening.}
\label{fig:core_inflation}
\source{Bank of Canada, Measures of Core Inflation database; Statistics Canada,
Table 18-10-0004-01. Values for 2025--2026 are estimates.}
\end{figure}

A critical insight from Figure~\ref{fig:core_inflation}: the gap between
headline and core inflation \textit{narrowed} in 2023--2024. As energy prices
fell back from their 2022 peaks, headline inflation declined sharply. But
core measures remained elevated well into 2024 because \term{services
inflation} --- wages, rents, and labour-intensive services --- proved
``sticky'' downward. This stickiness reflects the second-round effects of
inflation: when workers see prices rise, they demand higher wages; higher
wages increase firms' costs, pushing up prices further; and the cycle
reinforces itself. Breaking this cycle is the primary rationale for the
Bank of Canada's rate-hike strategy.

%==========================================================
\section{Real vs.\ Nominal Interest Rates: The Fisher Equation}
\label{sec:fisher_equation}
%==========================================================

When a bank offers you a savings account paying 5\% interest, is that a good
deal? The answer depends on inflation. If prices are rising at 3\% per year,
your 5\% nominal return only delivers 2\% more purchasing power --- the
real return is just 2\%. If inflation is running at 6\%, your ``5\%'' nominal
return is actually a loss of 1\% in purchasing power.

This insight is captured by the \term{Fisher equation}, one of the most
important relationships in monetary economics:

\begin{equation}
  r \approx i - \pi
  \label{eq:fisher_simple}
\end{equation}

\noindent where $r$ is the \term{real interest rate}, $i$ is the
\term{nominal interest rate}, and $\pi$ is the \term{inflation rate}.

The exact version is:

\begin{equation}
  (1 + r) = \frac{(1 + i)}{(1 + \pi)}
  \quad \Longleftrightarrow \quad
  r = \frac{1 + i}{1 + \pi} - 1
  \label{eq:fisher_exact}
\end{equation}

For the values typically encountered in Canada, the approximation
$r \approx i - \pi$ (Equation~\ref{eq:fisher_simple}) is sufficiently
accurate for analysis.

\begin{defbox}{Real vs.\ Nominal Interest Rates}
The \textbf{nominal interest rate} $i$ is the rate stated in the
loan contract or savings agreement --- the number on the Bank of Canada
website.

The \textbf{real interest rate} $r = i - \pi^e$ is the nominal rate
adjusted for (expected) inflation. It represents the true cost of
borrowing or the true return to saving in terms of purchasing power.

When making economic decisions, it is the \textit{real} interest rate
that matters. A firm deciding whether to invest in new equipment
compares the real borrowing cost against the real rate of return on
the equipment. A household deciding whether to take out a mortgage
compares the real cost of the loan against the real appreciation of
the home.

\textit{Ex ante vs.\ ex post:} Before the fact (\textit{ex ante}),
borrowers and lenders agree on a nominal rate based on their expectation
of future inflation $\pi^e$. After the fact (\textit{ex post}), the real
rate reflects actual inflation $\pi$. Unexpected inflation benefits
debtors (who repay in cheapened dollars) and hurts creditors (who receive
less purchasing power than expected).
\end{defbox}

\begin{table}[H]
\centering
\caption{Real Interest Rate Calculation, Canada: 2018--2026}
\label{tab:real_rates}
\renewcommand{\arraystretch}{1.30}
\begin{tabular}{lrrrrr}
\toprule
\textbf{Year} & \textbf{Nominal Rate} & \textbf{CPI Inflation} & \textbf{Real Rate}
              & \textbf{Real Rate} & \textbf{Policy} \\
              & $i$ (Bank Rate, \%) & $\pi$ (\%) & $r \approx i - \pi$ (\%)
              & Interpretation & Context \\
\midrule
2018 & 1.75 & 2.3 & $-$0.6 & Mildly stimulative & Late expansion \\
2019 & 1.75 & 1.9 & $-$0.2 & Roughly neutral    & Pre-COVID \\
2020 & 0.25 & 0.7 & $-$0.5 & Stimulative        & COVID emergency \\
2021 & 0.25 & 3.4 & $-$3.2 & Highly stimulative & Post-COVID stimulus \\
2022 & 4.25\rlap{$^*$} & 6.8 & $-$2.6\rlap{$^*$} & Still negative & Tightening cycle \\
2023 & 5.00 & 3.9 & +1.1   & Mildly restrictive & Tightening complete \\
2024 & 3.75\rlap{$^\dagger$} & 2.6 & +1.2 & Mildly restrictive & Easing begins \\
2025 & 3.00 & 2.1 & +0.9   & Near neutral       & Easing continues \\
2026 & 2.75 & 2.0 & +0.8   & Near neutral       & Soft landing \\
\bottomrule
\end{tabular}
\vspace{4pt}
\noindent\footnotesize{$^*$End-of-year Bank Rate; mid-year 2022 average rate was
approximately 2.0\%, making the real rate even more negative at $-4.8\%$.
$^\dagger$December 2024 rate after four cuts in the second half of 2024.}
\source{Bank of Canada key interest rates; Statistics Canada Table 18-10-0004-01;
author calculations.}
\end{table}

\noindent\textbf{The Key Lesson from Table~\ref{tab:real_rates}:}
In 2021 and through most of 2022, the real interest rate was deeply
\textit{negative} --- borrowers were effectively being paid to borrow in
real terms. At a mid-2022 Bank Rate of $i = 2.5\%$ and inflation of
$\pi = 7\%$, the real rate was approximately $r = -4.5\%$. This meant
that anyone borrowing at the prime rate to purchase real assets (housing,
stocks, commodities) was receiving a massive implicit subsidy from
inflation. Negative real rates were a powerful engine of the demand boom
that itself contributed to inflation.

%----------------------------------------------------------
\subsection{Real Wages: The Fisher Equation Applied to Labour}
%----------------------------------------------------------

The Fisher relationship applies equally to wages. The \term{nominal wage}
$W$ is the dollar amount a worker earns. The \term{real wage} is the nominal
wage deflated by the price level:

\begin{equation}
  w_{\text{real},t} = \frac{W_t}{\text{CPI}_t / 100}
  \label{eq:real_wage}
\end{equation}

The growth rate of real wages is approximately:

\begin{equation}
  g_{w,\text{real}} \approx g_{W,\text{nominal}} - \pi
  \label{eq:real_wage_growth}
\end{equation}

\noindent\textbf{Worked Example 2.2 --- Real Wage Growth, Canada 2021--2023}

Average hourly wages in Canada grew from \$28.57 in January 2021 to
\$32.36 in January 2023.
\begin{itemize}
  \item Nominal wage growth: $(32.36 - 28.57) / 28.57 \times 100 = 13.3\%$
  \item CPI in January 2021: 138.8; January 2023: 155.6
  \item CPI change: $(155.6 - 138.8)/138.8 \times 100 = 12.1\%$
  \item Real wage growth: $13.3\% - 12.1\% = +1.2\%$ over two years
\end{itemize}

Despite nominal wage gains that looked impressive, real wages grew by
only 0.6\% per year on average --- barely keeping pace with productivity.
For workers in the bottom two income quintiles, where wage gains were
smaller and inflation in food and shelter (which loom larger in their
budgets) was higher, real wages actually fell.

%==========================================================
\section{The Aggregate Demand--Aggregate Supply Framework}
\label{sec:adas}
%==========================================================

The \term{Aggregate Demand--Aggregate Supply (AD-AS) model} is the
workhorse framework for analysing macroeconomic fluctuations, including
inflation. It explains how the overall price level and real output are
jointly determined in the economy, and why they change over time.

%----------------------------------------------------------
\subsection{Aggregate Demand}
%----------------------------------------------------------

The \term{Aggregate Demand (AD) curve} shows the combinations of price
level and real output at which the goods market and money market are
simultaneously in equilibrium. It slopes downward for three reasons:

\begin{itemize}
  \item \textbf{Real balance effect:} A higher price level reduces the
        real value of household wealth (including cash and bonds held at
        face value), reducing consumption spending.
  \item \textbf{Interest rate effect:} A higher price level increases
        the demand for money (people need more nominal money to buy the
        same real quantities), pushing up interest rates and reducing
        investment.
  \item \textbf{Exchange rate effect:} Higher domestic prices make
        Canadian goods less competitive internationally, reducing exports
        and increasing imports --- reducing net exports $NX$.
\end{itemize}

A simplified aggregate demand relationship can be written:

\begin{equation}
  Y = \alpha_0 - \beta \cdot r + \gamma \cdot G - \delta \cdot P
  \label{eq:ad_simplified}
\end{equation}

\noindent where $Y$ is real output, $r$ is the real interest rate,
$G$ is government spending, and $P$ is the price level. The parameters
$\beta$ and $\delta$ capture the sensitivity of spending to interest
rates and prices, respectively.

\textbf{Shifts in AD:} The AD curve shifts rightward (increasing output
demand at every price level) when:
\begin{itemize}
  \item Government spending $G$ increases (fiscal expansion)
  \item The central bank cuts interest rates $r$ (monetary expansion)
  \item Household wealth or confidence rises
  \item Export demand increases (trading partner growth or currency depreciation)
\end{itemize}

%----------------------------------------------------------
\subsection{Aggregate Supply}
%----------------------------------------------------------

The \term{Short-Run Aggregate Supply (SRAS) curve} shows the total quantity
of goods and services firms are willing to produce at each price level,
given that input prices (especially wages) are fixed in the short run.
It slopes upward because higher output prices increase profit margins,
incentivizing firms to expand production.

A simple SRAS relationship:

\begin{equation}
  P = P^e + \frac{1}{\alpha}(Y - Y_f)
  \label{eq:sras}
\end{equation}

\noindent where $P$ is the current price level, $P^e$ is the expected
price level (determining the position of the SRAS curve), $Y_f$ is
potential output (full-employment output), and $\alpha$ is a positive
parameter governing how responsive prices are to output gaps.

The \term{Long-Run Aggregate Supply (LRAS) curve} is vertical at
$Y = Y_f$ --- potential output. In the long run, all prices are flexible,
the economy returns to $Y_f$, and the price level adjusts to whatever is
required for this equilibrium. Monetary policy cannot permanently raise
output above $Y_f$; it can only affect the price level in the long run.

%----------------------------------------------------------
\subsection{Two Types of Inflation}
%----------------------------------------------------------

\begin{figure}[H]
\centering
\begin{tikzpicture}[
  node distance=0cm,
  axisline/.style={-Stealth, thick, color=gray!70},
  curvep/.style={thick, color=policyblue},
  curvea/.style={thick, color=casered},
  curved/.style={thick, color=unbcgreen},
  curveas/.style={thick, color=dataorange},
  labelstyle/.style={font=\small}
]

% ---- PANEL A: DEMAND-PULL ----
\begin{scope}[xshift=0cm]
\node[font=\small\bfseries] at (3.2, 5.3) {Panel A: Demand-Pull Inflation};
% Axes
\draw[axisline] (0,0) -- (0,5) node[above, font=\small] {$P$};
\draw[axisline] (0,0) -- (6.2,0) node[right, font=\small] {$Y$};
% LRAS vertical
\draw[dashed, gray, thick] (3.5,0) -- (3.5,4.8)
  node[above, font=\footnotesize] {LRAS};
% SRAS curve (upward sloping)
\draw[curveas, thick] (0.5,0.5) .. controls (2,1.8) .. (5.5,4.5)
  node[right, font=\footnotesize, color=dataorange] {SRAS};
% AD1 (original)
\draw[curved] (0.5,4.2) .. controls (2.5,2.5) .. (5.5,0.5)
  node[right, font=\footnotesize, color=unbcgreen] {AD$_1$};
% AD2 (rightward shift)
\draw[curvep, dashed] (1.5,4.5) .. controls (3.5,2.8) .. (6.0,0.8)
  node[right, font=\footnotesize, color=policyblue] {AD$_2$};
% Equilibria
\fill[unbcgreen] (3.0,2.0) circle (3pt);
\fill[policyblue] (4.2,2.9) circle (3pt);
% Labels
\draw[dashed, gray!50, thin] (3.0,0) -- (3.0,2.0) -- (0,2.0);
\draw[dashed, gray!50, thin] (4.2,0) -- (4.2,2.9) -- (0,2.9);
\node[left, font=\footnotesize] at (0,2.0) {$P_1$};
\node[left, font=\footnotesize] at (0,2.9) {$P_2$};
\node[below, font=\footnotesize] at (3.0,0) {$Y_1$};
\node[below, font=\footnotesize] at (4.2,0) {$Y_2$};
\node[below, font=\footnotesize] at (3.5,-0.05) {$Y_f$};
% Arrow showing shift
\draw[->, thick, policyblue!60] (2.8,2.6) -- (3.6,2.6);
\end{scope}

% ---- PANEL B: COST-PUSH ----
\begin{scope}[xshift=7.5cm]
\node[font=\small\bfseries] at (3.2, 5.3) {Panel B: Cost-Push Inflation};
% Axes
\draw[axisline] (0,0) -- (0,5) node[above, font=\small] {$P$};
\draw[axisline] (0,0) -- (6.2,0) node[right, font=\small] {$Y$};
% LRAS vertical
\draw[dashed, gray, thick] (3.5,0) -- (3.5,4.8)
  node[above, font=\footnotesize] {LRAS};
% SRAS1 (original)
\draw[curveas] (0.5,0.5) .. controls (2,1.8) .. (5.5,4.5)
  node[right, font=\footnotesize, color=dataorange] {SRAS$_1$};
% SRAS2 (leftward shift --- cost shock)
\draw[casered, dashed, thick] (0.2,1.5) .. controls (1.5,2.8) .. (4.8,5.0)
  node[right, font=\footnotesize, color=casered] {SRAS$_2$};
% AD curve
\draw[curved] (0.5,4.2) .. controls (2.5,2.5) .. (5.5,0.5)
  node[right, font=\footnotesize, color=unbcgreen] {AD};
% Equilibria
\fill[unbcgreen] (3.0,2.0) circle (3pt);
\fill[casered] (2.1,2.9) circle (3pt);
% Labels
\draw[dashed, gray!50, thin] (3.0,0) -- (3.0,2.0) -- (0,2.0);
\draw[dashed, gray!50, thin] (2.1,0) -- (2.1,2.9) -- (0,2.9);
\node[left, font=\footnotesize] at (0,2.0) {$P_1$};
\node[left, font=\footnotesize] at (0,2.9) {$P_2$};
\node[below, font=\footnotesize] at (3.0,0) {$Y_1=Y_f$};
\node[below, font=\footnotesize] at (2.1,0) {$Y_2$};
% Arrow showing leftward SRAS shift
\draw[->, thick, casered!70] (2.0,2.3) -- (1.2,2.3);
% Stagflation label
\node[font=\footnotesize, casered] at (3.5,4.0)
  {\textit{Stagflation}};
\node[font=\footnotesize, casered] at (3.5,3.6)
  {($P\uparrow$, $Y\downarrow$)};
\end{scope}

\end{tikzpicture}
\caption{Demand-Pull vs.\ Cost-Push Inflation in the AD-AS Framework.
\textbf{Panel A (Demand-Pull):} A rightward shift of AD (from AD$_1$ to
AD$_2$) --- driven by fiscal stimulus, low interest rates, or rising
household wealth --- moves the economy from equilibrium $(Y_1, P_1)$ to
$(Y_2, P_2)$, raising both output and prices. If $Y_2 > Y_f$, the economy
is overheating. \textbf{Panel B (Cost-Push):} A leftward shift of SRAS ---
driven by higher energy prices, supply chain disruptions, or rising input
costs --- produces the combination of higher prices \textit{and} lower
output known as \textbf{stagflation}. Canada's 2021--2023 inflation episode
involved both mechanisms simultaneously: massive AD expansion from COVID
stimulus and near-zero rates, combined with supply chain bottlenecks and
the Russian-invasion energy shock.}
\label{fig:adas}
\end{figure}

\begin{canexample}{Canada's 2021--2023 Inflation: Both Demand and Supply}
Canada's inflation surge was unusual because both AD and AS shocks
occurred simultaneously and reinforced each other.

\textbf{Demand side:} The federal government injected approximately
\$350 billion in emergency spending between 2020 and 2022. The Bank of
Canada held its policy rate at 0.25\% until March 2022. Household savings
rates surged during lockdowns (\$250+ billion in excess savings accumulated),
creating pent-up demand that exploded when the economy reopened. These
forces shifted AD rightward continuously from mid-2020 through 2021.

\textbf{Supply side:} Global shipping container costs rose tenfold.
Semiconductor shortages idled automotive production lines --- used car
prices in Canada rose over 40\% in 2021. Russia's invasion of Ukraine
(February 2022) triggered an energy and food commodity shock, shifting
SRAS leftward. Labour shortages in food processing, transportation, and
construction added further upward cost pressure.

\textbf{Result:} In the AD-AS diagram, Canada experienced both a large
rightward AD shift and a leftward SRAS shift. Prices rose sharply. Output
expanded beyond potential in 2021 (generating an ``output gap'' of
approximately +2\% of potential GDP) before the rate hikes began
cooling demand. This combination --- demand exceeding supply on both
sides --- is precisely what generates the most acute inflationary episodes.
\end{canexample}

%==========================================================
\section{The Phillips Curve: Inflation and Unemployment}
\label{sec:phillips_curve}
%==========================================================

The \term{Phillips Curve} describes the empirical and theoretical
relationship between inflation and unemployment. The original version,
identified by economist A.W. Phillips (1958) using UK data from 1861--1957,
showed a stable negative relationship: when unemployment was low, wages
(and thus prices) tended to rise faster.

%----------------------------------------------------------
\subsection{The Short-Run Phillips Curve}
%----------------------------------------------------------

The short-run Phillips Curve can be expressed as:

\begin{equation}
  \pi_t = \pi^e_t - \alpha(u_t - u_n)
  \label{eq:phillips_curve}
\end{equation}

\noindent where:
\begin{itemize}
  \item $\pi_t$ is the actual inflation rate in period $t$
  \item $\pi^e_t$ is the expected inflation rate (which shifts the curve)
  \item $u_t$ is the actual unemployment rate
  \item $u_n$ is the \term{natural rate of unemployment} (the NAIRU ---
        Non-Accelerating Inflation Rate of Unemployment)
  \item $\alpha > 0$ is a parameter measuring how sensitive inflation is
        to labour market slack
\end{itemize}

The term $(u_t - u_n)$ is the \term{unemployment gap}. When $u_t < u_n$
(labour markets are tight, unemployment below natural rate), there is
upward pressure on wages and prices: $\pi_t > \pi^e_t$. When
$u_t > u_n$ (slack in labour markets), inflation tends to run below
expectations: $\pi_t < \pi^e_t$.

For Canada, the Bank of Canada estimates the natural rate of unemployment
$u_n$ at approximately 5.5--6.0\% in the mid-2020s (accounting for
structural factors such as the geographic mismatch between job vacancies
and unemployed workers, and skills mismatches). When unemployment fell
to 4.9\% in mid-2022, the gap $(u_t - u_n) \approx -1.0$ to $-1.5$
percentage points, contributing upward pressure on wages and prices.

%----------------------------------------------------------
\subsection{The Long-Run Phillips Curve}
%----------------------------------------------------------

If workers have \term{rational expectations} --- they correctly anticipate
inflation on average --- then any attempt by the central bank to exploit
the inflation--unemployment trade-off will be defeated. If the Bank
tolerates higher inflation to keep unemployment below the natural rate,
workers eventually demand higher nominal wages to compensate for the
higher expected inflation, shifting the SRAS curve upward and returning
the economy to $u_n$ --- but now at a permanently higher inflation rate.

The long-run Phillips Curve is therefore \textbf{vertical at $u = u_n$}:
there is no long-run trade-off between inflation and unemployment. This
result, developed independently by Edmund Phelps and Milton Friedman in
1968, is one of the most important insights in macroeconomics. It implies:
\begin{itemize}
  \item Monetary policy cannot permanently lower unemployment below $u_n$.
  \item Attempting to do so only generates accelerating inflation.
  \item The only permanent gains in employment come from structural
        policies that reduce $u_n$ itself (better job matching, education,
        reduced geographical barriers to labour mobility).
\end{itemize}

%----------------------------------------------------------
\subsection{The Sacrifice Ratio}
%----------------------------------------------------------

When a central bank wants to \textit{reduce} inflation that has become
embedded in expectations, it must push unemployment above the natural
rate --- accepting a period of higher unemployment. The \term{sacrifice
ratio} measures the cost of this disinflation:

\begin{equation}
  \text{Sacrifice Ratio} = \frac{\text{Cumulative excess unemployment (\%)}}
                                 {\text{Reduction in inflation (pp)}}
  \label{eq:sacrifice_ratio}
\end{equation}

Empirical estimates for Canada typically place the sacrifice ratio between
1.5 and 3.0: to permanently reduce inflation by 1 percentage point requires
approximately 1.5--3.0 percentage points of excess unemployment sustained
for one year. For the 2022--2024 disinflation of approximately 4 percentage
points (from 6.8\% to 2.6\%), the implied employment cost was significant.

\begin{figure}[H]
\centering
\begin{tikzpicture}
\begin{axis}[
  width=10cm, height=8cm,
  xlabel={Unemployment Rate (\%)},
  ylabel={CPI Inflation, Year-over-Year (\%)},
  xlabel style={font=\small},
  ylabel style={font=\small},
  xmin=4.5, xmax=10.5,
  ymin=-0.5, ymax=8.5,
  xtick={5,6,7,8,9,10},
  ytick={0,1,2,3,4,5,6,7,8},
  xticklabel style={font=\small},
  yticklabel style={font=\small},
  ymajorgrids=true, xmajorgrids=true,
  grid style={dashed, gray!20},
  clip=false
]
% 2% inflation target line
\addplot[unbcgreen!50, dashed, thin] coordinates {(4.5,2)(10.5,2)};
% Natural rate reference
\addplot[gray!50, dotted, thin] coordinates {(5.8,0)(5.8,8.5)}
  node[above, font=\footnotesize, text=gray] {$u_n \approx 5.8\%$};
% Data points with year labels
\addplot[only marks, mark=*, mark size=3pt, policyblue]
  coordinates {
    (6.9,1.1)(7.0,1.4)(6.3,1.6)(5.8,2.3)(5.7,1.9)
    (9.5,0.7)(7.5,3.4)(5.3,6.8)(5.4,3.9)(6.3,2.6)(6.7,2.1)(6.5,2.0)
  };
% Year labels
\node[font=\footnotesize, policyblue, above right] at (axis cs:6.9,1.1) {15};
\node[font=\footnotesize, policyblue, above right] at (axis cs:7.0,1.4) {16};
\node[font=\footnotesize, policyblue, above right] at (axis cs:6.3,1.6) {17};
\node[font=\footnotesize, policyblue, above right] at (axis cs:5.8,2.3) {18};
\node[font=\footnotesize, policyblue, above right] at (axis cs:5.7,1.9) {19};
\node[font=\footnotesize, policyblue, right] at (axis cs:9.5,0.7) {20};
\node[font=\footnotesize, policyblue, above right] at (axis cs:7.5,3.4) {21};
\node[font=\footnotesize, policyblue, above] at (axis cs:5.3,6.8) {22};
\node[font=\footnotesize, policyblue, right] at (axis cs:5.4,3.9) {23};
\node[font=\footnotesize, policyblue, below right] at (axis cs:6.3,2.6) {24};
\node[font=\footnotesize, policyblue, below right] at (axis cs:6.7,2.1) {25};
\node[font=\footnotesize, policyblue, below right] at (axis cs:6.5,2.0) {26};
% Arrow showing 2022 surge direction
\draw[->, casered, thick] (axis cs:7.5,3.4) -- (axis cs:5.4,6.5);
\node[font=\footnotesize, casered, rotate=45] at (axis cs:6.4,5.2) {rate surge};
% Arrow showing disinflation
\draw[->, unbcgreen!80, thick] (axis cs:5.4,6.7) to[bend right=15] (axis cs:6.7,2.2);
\node[font=\footnotesize, unbcgreen!80] at (axis cs:4.9,4.5) {disinflation};
\end{axis}
\end{tikzpicture}
\caption{Phillips Curve: Canada, 2015--2026. Each point plots the annual
unemployment rate against the annual average CPI inflation rate, labelled
by the last two digits of the year. The broad negative relationship
--- lower unemployment associated with higher inflation --- is visible
in 2015--2022, consistent with the short-run Phillips Curve. The 2020
point (year 20) is an outlier: unemployment spiked to 9.5\% while
inflation collapsed to 0.7\%, reflecting the demand destruction of
the COVID lockdowns. The 2022 point (year 22) sits at the extreme
top-left: very low unemployment (5.3\%) and very high inflation (6.8\%),
reflecting a severely overheated economy. The red arrow traces the
surge and the green arrow traces the subsequent disinflation path,
along which unemployment rose as inflation fell.}
\label{fig:phillips_curve}
\source{Statistics Canada, Tables 14-10-0287-01 and 18-10-0004-01.
Annual averages; 2025--2026 estimated.}
\end{figure}

%==========================================================
\section{Canada's Inflation-Targeting Framework}
\label{sec:inflation_targeting}
%==========================================================

Canada was the second country in the world (after New Zealand in 1990)
to adopt a formal \term{inflation target}. The framework was introduced
in February 1991 under a joint agreement between the Bank of Canada and
the federal government. It has been renewed multiple times, most recently
in 2021 with modifications that reflected the lessons of the COVID era.

The key features of the current framework:

\begin{itemize}
  \item \textbf{Target:} 2\% measured CPI inflation.
  \item \textbf{Control range:} 1--3\% (the Bank uses the range as an
        indicator of acceptable outcomes, not as hard bounds).
  \item \textbf{Horizon:} 18--24 months --- the time over which the Bank
        aims to return inflation to the 2\% target following a shock.
        This reflects the lags in monetary policy transmission.
  \item \textbf{Instrument:} The overnight rate (policy rate).
  \item \textbf{Communication:} Eight pre-announced Monetary Policy
        Announcement (MPA) dates per year, accompanied by a quarterly
        Monetary Policy Report (MPR) with detailed economic forecasts.
  \item \textbf{Accountability:} The Bank Governor must write an open
        letter to the Minister of Finance explaining why inflation has
        moved more than 1 percentage point above or below the 2\% target.
\end{itemize}

\textbf{Why targeting works --- the role of credibility:}
Inflation targeting is more than a rule; it is a commitment device.
When firms, workers, and financial markets believe the Bank will keep
inflation at 2\%, they set wages and prices consistent with 2\% inflation.
This self-fulfilling property --- ``credibility'' --- makes disinflation
less costly (smaller sacrifice ratio) and makes it easier to keep
inflation at target even when shocks occur. If the Bank loses credibility,
inflation expectations become de-anchored, requiring more aggressive
tightening to achieve the same disinflationary effect.

The 2021--2022 episode tested this credibility severely. By maintaining
the policy rate at 0.25\% while inflation surpassed 5\% --- citing
concerns about premature withdrawal of pandemic support --- the Bank
allowed long-term inflation expectations to creep above 2\%. By the time
rates rose aggressively, some of the benefit of credibility had been
partially eroded, requiring more tightening than if the Bank had acted
earlier.

%==========================================================
\section{Monetary Policy Transmission}
\label{sec:transmission}
%==========================================================

When the Bank of Canada changes its overnight rate, the effects on inflation
and output do not occur immediately. There are multiple channels through
which the policy rate change propagates through the economy, and each
channel operates with different lags. The full effect on inflation is
typically estimated to take \textbf{18--24 months}.

\begin{defbox}{The Four Channels of Monetary Policy Transmission}
\textbf{1. Credit Cost Channel:} Higher overnight rates raise the cost
of borrowing for households and firms. Mortgage rates, car loan rates,
business credit lines, and corporate bond yields all rise. This reduces
investment spending (firms borrow less to expand) and consumer spending
(households reduce debt-financed purchases). This is the fastest and
most direct channel.

\textbf{2. Asset Price Channel:} Higher interest rates reduce the
present value of future cash flows, lowering stock prices, bond prices,
and real estate valuations. Lower household wealth reduces consumption
through the \term{wealth effect}. A fall in equity prices also raises
firms' cost of capital, further reducing investment.

\textbf{3. Exchange Rate Channel:} Higher Canadian interest rates attract
foreign capital seeking superior returns, increasing demand for Canadian
dollars and appreciating the exchange rate. A stronger Canadian dollar
makes imports cheaper (reducing the domestic price of imported goods)
and makes Canadian exports less competitive internationally (reducing
net exports). Both effects reduce domestic price pressure.

\textbf{4. Expectations Channel:} When the Bank credibly signals
its inflation intentions, households and firms adjust their price- and
wage-setting behaviour immediately, \textit{before} the other channels
fully operate. If workers and firms believe the Bank will achieve 2\%
inflation, they set wages and prices consistent with 2\% --- effectively
doing the Bank's job for it. This is why central bank communication
(\term{forward guidance}) is itself a monetary policy instrument.
\end{defbox}

The relative importance of these channels in Canada reflects the
structure of the Canadian economy. The credit cost channel is particularly
powerful because:

\begin{enumerate}
  \item Canada has one of the highest household debt-to-income ratios
        among OECD countries (~180\% as of 2023), making households
        extremely sensitive to interest rate changes.
  \item Approximately 35\% of Canadian mortgages were in variable-rate
        products as of 2022, meaning rate increases transmitted to
        mortgage payments with minimal lag.
  \item Housing wealth constitutes approximately 40\% of total Canadian
        household net worth, making the asset price channel powerful.
\end{enumerate}

%==========================================================
\section{Case Studies}
\label{sec:ch2_casestudies}
%==========================================================

%----------------------------------------------------------
\begin{casestudy}{The Bank of Canada's 2022--2024 Rate-Hike Cycle}

\textbf{Background:} By early 2022, Canada's inflation had risen to levels
not seen since the early 1980s. The Bank of Canada, which had kept its
policy rate at the emergency floor of 0.25\% since March 2020, faced the
challenge of rapidly tightening monetary policy without triggering a
severe recession.

\textbf{The Timeline:} The tightening cycle began on March 2, 2022, with
a 25-basis-point increase --- the Bank's first rate hike since October 2018.
What followed was the most rapid tightening cycle in the Bank's modern
history.

\begin{table}[H]
\centering
\caption{Bank of Canada Rate Decisions, March 2022 -- December 2024}
\label{tab:rate_decisions}
\renewcommand{\arraystretch}{1.25}
\begin{tabular}{lrrl}
\toprule
\textbf{Date} & \textbf{Change (bps)} & \textbf{New Rate (\%)}
              & \textbf{Bank's Primary Rationale} \\
\midrule
Mar 2, 2022  & +25  & 0.50 & Inflation exceeding target; recovery solid \\
Apr 13, 2022 & +50  & 1.00 & Inflation broad-based; front-loading needed \\
Jun 1, 2022  & +50  & 1.50 & CPI 7.7\% in Apr; labour markets tight \\
Jul 13, 2022 & +100 & 2.50 & CPI 8.1\% June peak; unconventional move \\
Sep 7, 2022  & +75  & 3.25 & Core measures still rising \\
Oct 26, 2022 & +50  & 3.75 & Moderate; economy showing some cooling \\
Dec 7, 2022  & +50  & 4.25 & Inflation still above target \\
Jan 25, 2023 & +25  & 4.50 & Conditional pause signalled \\
Mar 8, 2023  & +25  & 4.75 & Core inflation re-accelerated \\
Jun 7, 2023  & +25  & 5.00 & GDP and CPI stronger than expected \\
\midrule
\textit{Jul--Dec 2023} & \textit{Held} & \textit{5.00} & \textit{Assessment period} \\
\midrule
Jun 5, 2024  & $-$25 & 4.75 & Inflation returning to target \\
Jul 24, 2024 & $-$25 & 4.50 & Labour market easing \\
Sep 4, 2024  & $-$25 & 4.25 & CPI at 2.5\%; output gap opened \\
Oct 23, 2024 & $-$50 & 3.75 & Downside risks growing \\
Dec 11, 2024 & $-$50 & 3.25 & Inflation at target; growth below potential \\
\bottomrule
\end{tabular}
\source{Bank of Canada, Monetary Policy Announcements, 2022--2024.
``bps'' = basis points (1 basis point = 0.01 percentage point).}
\end{table}

\textbf{The July 2022 Anomaly:} The July 13, 2022 decision --- a 100-basis-point
increase, the largest single-meeting rate hike since 1998 --- was remarkable.
The Bank cited the June 2022 CPI print of 8.1\% and the risk that
inflation expectations would become ``de-anchored'' from the 2\% target.
At that point, the Bank was willing to accept a sharp slowdown rather than
allow inflation to become entrenched.

\textbf{Outcomes:} By June 2024, inflation had returned to 2.7\% --- near
the 2\% target. Canada avoided a technical recession (two consecutive
quarters of negative real GDP growth) in 2022--2023, though real GDP per
capita did decline. The unemployment rate rose from a 50-year low of 4.9\%
in mid-2022 to 6.5--6.8\% by late 2024.

\textbf{Evaluation:} Did the Bank of Canada succeed? On price stability,
yes: inflation returned to target within two years of the June 2022 peak.
On the real economy, the outcome was relatively benign by historical
standards --- no deep recession, modest unemployment increase, and nominal
wage growth that largely kept pace with declining inflation. However,
mortgage holders who renewed into 5\% rates experienced the most severe
payment shock in a generation, and some economists argue the Bank held
rates too high for too long in 2023--24, unnecessarily slowing growth.
\end{casestudy}

%----------------------------------------------------------
\begin{casestudy}{Shelter Inflation: The CPI's Most Powerful Component}

\textbf{The Issue:} Table~\ref{tab:cpi_components} shows that shelter
inflation remained elevated --- and in some measures \textit{accelerating}
--- into 2023 and 2024, even as energy and goods prices eased. Mortgage
interest costs rose 28.5\% in 2023 and 19.4\% in 2024 year-over-year,
making shelter the single largest driver of measured inflation in the
post-peak period.

\textbf{The Paradox:} The Bank of Canada raised interest rates to fight
inflation. But because mortgage interest costs are included in the CPI
shelter component (they are part of the cost of ``owning'' a home),
the rate hikes directly increased the CPI they were supposed to reduce.
This creates a short-run paradox: tighter monetary policy mechanically
raises measured inflation through the mortgage channel, even as it slows
the economy and reduces demand-side inflation pressure.

\textbf{Why Does This Happen in Canada?}
Canada's CPI shelter methodology, unlike that of many other countries,
includes mortgage interest cost (MIC) directly. The US CPI, by contrast,
uses \term{owners' equivalent rent} --- an estimate of what homeowners
would pay if they rented their homes --- which is more insulated from
short-run interest rate movements.

\textbf{Policy Implication:} This is one reason the Bank of Canada
monitors core inflation measures (particularly CPI-trim and CPI-median)
that tend to dampen the direct mortgage cost effect. It is also why
economists sometimes argue that the headline CPI in Canada gives a
misleading picture of ``true'' cost-of-living inflation during monetary
tightening cycles.

\textbf{Quantitative impact:} In 2023, mortgage interest costs contributed
approximately +1.4 percentage points to Canada's headline inflation of 3.9\%
--- meaning that without the mortgage component, headline inflation would
have been closer to 2.5\%. The same rate hikes that were cooling inflation
were adding roughly 1.4 points to the inflation number, creating both
an economic paradox and a political communications challenge for the Bank.
\end{casestudy}

%----------------------------------------------------------
\begin{casestudy}{Who Pays for Inflation? Distributional Consequences Across Income Groups}

\textbf{Background:} The CPI measures inflation for the ``average'' household.
But average households do not exist --- spending patterns, asset holdings,
debt levels, and income sources vary enormously across the income distribution.
The 2022 inflation surge was not experienced equally.

\textbf{The Regressivity of Inflation:} Low-income households devote larger
shares of their budgets to food, shelter, and energy --- the three categories
that rose fastest in 2022. Statistics Canada analysis estimates that the
effective inflation rate experienced by the bottom income quintile in
2022 was approximately 8.5\%, compared to approximately 6.8\% for the top
quintile. A gap of 1.7 percentage points may seem small in absolute terms,
but for a household spending 95\% of its income, the difference between
an 8.5\% and 6.8\% inflation rate is the difference between severe hardship
and moderate difficulty.

\textbf{Debtors and Creditors:} Unexpected inflation redistributes wealth
from creditors to debtors. Anyone holding a fixed-rate loan from before
the inflation surge found that the real value of their debt was being
eroded by inflation. Mortgage holders with 5-year fixed rates locked in
at 2\% in 2020 enjoyed negative real borrowing costs throughout 2022 ---
their 2\% nominal rate minus 6\%+ inflation implied a real borrowing cost
of approximately $-4\%$. Retirees and pension funds holding fixed-income
securities saw the real value of their assets eroded.

\textbf{Wage-earners vs.\ Capital owners:} Workers whose wages did not
fully adjust to inflation experienced real pay cuts. Workers in unionized
sectors with cost-of-living adjustment (COLA) clauses in their contracts
were partially protected. Non-unionized, part-time, and gig workers ---
disproportionately low-income, racialized, and young --- had the least
wage protection. Capital owners who held real assets (housing, stocks,
commodities) saw their nominal values rise with inflation, preserving
real wealth while nominal wealth grew.

\textbf{Policy Implications:} The distributional consequences of inflation
strengthen the case for price stability as a distributive goal, not just
an efficiency goal. Inflation is not a politically neutral phenomenon:
it systematically favours debtors over creditors, real asset holders
over cash savers, and the flexible-wage over the fixed-wage workforce.
\end{casestudy}

%==========================================================
\section{Distributional Effects of Inflation}
\label{sec:distributional}
%==========================================================

The case study above is worth formalising. Figure~\ref{fig:inflation_burden}
presents estimated effective inflation rates by income quintile for Canada
in 2022, constructed by reweighting CPI components using the expenditure
patterns of each quintile from the Survey of Household Spending.

\begin{figure}[H]
\centering
\begin{tikzpicture}
\begin{axis}[
  ybar,
  bar width=0.8cm,
  width=10cm, height=6.5cm,
  symbolic x coords={{Q1 (Lowest)},{Q2},{Q3},{Q4},{Q5 (Highest)}},
  xtick=data,
  x tick label style={font=\small, rotate=20, anchor=east},
  ylabel={Effective Inflation Rate (\%), 2022},
  ylabel style={font=\small},
  ymin=5.5, ymax=9.5,
  ytick={6.0,6.5,7.0,7.5,8.0,8.5,9.0},
  yticklabel style={font=\small},
  ymajorgrids=true,
  grid style={dashed, gray!25},
  enlarge x limits=0.15,
  nodes near coords,
  nodes near coords style={font=\small\bfseries},
  every axis plot/.append style={fill=casered!65, draw=casered!85}
]
\addplot coordinates {
  ({Q1 (Lowest)}, 8.5)
  ({Q2},          7.8)
  ({Q3},          7.5)
  ({Q4},          7.2)
  ({Q5 (Highest)},6.8)
};
% Reference line for official headline CPI
\addplot[unbcgreen, dashed, thick, forget plot] coordinates
  {({Q1 (Lowest)},6.8)({Q5 (Highest)},6.8)}
  node[right, font=\footnotesize, color=unbcgreen, xshift=2pt]
  {Headline CPI (6.8\%)};
\end{axis}
\end{tikzpicture}
\caption{Estimated Effective Inflation Rate by Income Quintile, Canada, 2022.
Low-income households (Q1) experienced approximately 8.5\% effective inflation
--- 1.7 percentage points above the headline CPI of 6.8\% --- because
they spend a larger share of their budgets on food, shelter, and energy,
the three categories with the sharpest price increases in 2022. The
dashed line shows the official headline CPI for reference. Inflation
is not a neutral tax: it falls disproportionately on those least able to
absorb it.}
\label{fig:inflation_burden}
\source{Author calculations using Statistics Canada Survey of Household
Spending and CPI component data. Methodology follows Argente and Lee (2021)
and Statistics Canada analytical studies.}
\end{figure}

%==========================================================
\section{The Rate-Hike Cycle: Nominal and Real Rates}
\label{sec:rate_cycle}
%==========================================================

\begin{figure}[H]
\centering
\begin{tikzpicture}
\begin{axis}[
  width=14cm, height=7cm,
  xlabel={Year},
  ylabel={Interest Rate / Inflation Rate (\%)},
  xlabel style={font=\small},
  ylabel style={font=\small},
  xmin=2017.8, xmax=2026.5,
  ymin=-6, ymax=7,
  xtick={2018,2019,2020,2021,2022,2023,2024,2025,2026},
  xticklabel style={font=\small},
  ytick={-6,-4,-2,0,2,4,6},
  yticklabel style={font=\small},
  ymajorgrids=true,
  grid style={dashed, gray!20},
  legend pos=north west,
  legend style={font=\small},
  clip=false
]
% Zero line
\draw[gray!50, thin] (axis cs:2017.8,0) -- (axis cs:2026.5,0);
% Neutral real rate reference
\addplot[gray!40, dotted] coordinates {(2017.8,1.0)(2026.5,1.0)}
  node[right, font=\footnotesize, color=gray!60] {$r^* \approx 1\%$};
% Overnight rate (nominal)
\addplot[policyblue, thick, mark=square*, mark size=2pt] coordinates {
  (2018,1.75)(2019,1.75)(2020,0.25)(2021,0.25)
  (2022,4.25)(2023,5.00)(2024,3.25)(2025,3.00)(2026,2.75)
};
\addlegendentry{Nominal overnight rate $i$}
% Real interest rate (i - pi)
\addplot[casered, thick, mark=triangle*, mark size=2.5pt] coordinates {
  (2018,-0.55)(2019,-0.15)(2020,-0.45)(2021,-3.15)
  (2022,-2.55)(2023,1.10)(2024,0.65)(2025,0.90)(2026,0.75)
};
\addlegendentry{Real rate $r = i - \pi$}
% Shaded negative real rate area
\fill[casered!10] (axis cs:2017.8,0) -- (axis cs:2018,-0.55)
  -- (axis cs:2019,-0.15) -- (axis cs:2020,-0.45) -- (axis cs:2021,-3.15)
  -- (axis cs:2022,-2.55) -- (axis cs:2022.8,0) -- cycle;
\node[font=\footnotesize, casered, rotate=0] at (axis cs:2020.5,-2.0)
  {Negative real rates};
\node[font=\footnotesize, casered] at (axis cs:2021,-3.5)
  {(fuel for inflation)};
\end{axis}
\end{tikzpicture}
\caption{Nominal Overnight Rate vs.\ Real Interest Rate, Canada, 2018--2026.
The shaded region marks the period of negative real interest rates (2018--2022).
In 2021, the real rate fell to approximately $-3.2\%$ --- deeply negative ---
as the nominal rate was held at 0.25\% while inflation rose to 3.4\%.
By mid-2022, the combination of rapid nominal rate increases and peak inflation
meant the real rate was still close to $-4.5\%$ even as nominal rates rose
to 2.5\%. Only in 2023, when nominal rates reached 5.00\% and inflation fell
to 3.9\%, did the real rate finally turn positive. The return to positive
real rates was the mechanism through which monetary policy finally cooled
aggregate demand.}
\label{fig:real_nominal_rates}
\source{Bank of Canada key interest rates; Statistics Canada Table 18-10-0004-01;
author calculations using Fisher approximation $r \approx i - \pi$.}
\end{figure}

%==========================================================
\section*{Policy Debate}
%==========================================================

\begin{policydebate}{Did the Bank of Canada Hike Too Far, Too Fast --- or Not Far Enough, Too Late?}

\textbf{The Issue:} The Bank of Canada's 2022--2024 rate-hike cycle was
the most aggressive in its modern history. Whether it was the right response
depends critically on the diagnosis of what caused the inflation and on
judgements about the costs of different types of policy error.

\textbf{Case 1 --- Too Slow to Act (The Hawkish Critique):}
Critics argue the Bank should have raised rates earlier --- potentially
as early as the summer of 2021, when CPI was already at 3.7\%. By
maintaining emergency-level rates through 2021, the Bank:
\begin{itemize}
  \item Allowed real interest rates to fall to $-3\%$, massively stimulating
        borrowing and spending.
  \item Failed to ``lean against'' surging housing prices, which added
        to the wealth effect and the inflation spiral.
  \item Let inflation expectations drift above 2\%, requiring more
        aggressive tightening later.
\end{itemize}
The implicit sacrifice ratio of this delayed response was higher than
it needed to be: more unemployment was ultimately required to achieve
the same disinflation because expectations had partially de-anchored.

\textbf{Case 2 --- Too Aggressive (The Dovish Critique):}
Others argue the Bank moved too fast and too far, particularly after
mid-2022. By the time the July 2022 100-basis-point hike was made,
global supply chains were already improving, energy prices were
beginning to fall, and the inflation peak was likely already past.
The aggressive tightening:
\begin{itemize}
  \item Imposed mortgage payment shocks on millions of Canadians who had
        no role in causing inflation.
  \item Slowed real GDP per capita growth below zero for five consecutive
        quarters in 2023--24.
  \item Risked pushing unemployment unnecessarily high, creating real
        human costs (job loss, bankruptcy, reduced business investment)
        to achieve a disinflation that was partly happening on its own
        as supply chains normalized.
\end{itemize}

\textbf{The Evidence:}
\begin{itemize}
  \item Canada's peak inflation (8.1\%) was lower than the UK (11.1\%)
        and comparable to the US (9.1\%).
  \item The Bank achieved its inflation target by mid-2024 without a
        technical recession.
  \item Real GDP per capita nonetheless fell for several quarters ---
        a ``per capita recession'' even without the formal definition
        being triggered.
  \item Variable-rate mortgage holders saw monthly payments increase
        by an average of approximately \$700--\$900.
  \item Business insolvencies rose to near-record levels in 2024.
\end{itemize}

\textbf{The Analytical Framework:}
Evaluating central bank decisions requires distinguishing what the Bank
\textit{knew} at the time from what we know now. Using today's data to
critique 2021--2022 decisions is ``Monday morning quarterbacking.'' The
correct standard is: given the information available in real time, were
the Bank's decisions within the range of reasonable policy choices?
Most economists conclude the answer is yes --- though many would have
preferred earlier action in 2021 and possibly a slightly lower terminal
rate in 2023. This is a case where the difference between a good outcome
and a great outcome was perhaps 50--75 basis points of peak rate ---
a reminder of how consequential seemingly small numerical differences
in monetary policy can be.

\textbf{Your Position:} After reading this chapter, where do you stand?
Use the Fisher equation, the AS-AD model, and the Phillips Curve to
structure your argument.
\end{policydebate}

%==========================================================
\begin{everydaylife}{How Inflation Affects Your Grocery Bill --- and What You Can Do About It}

When Statistics Canada reports that food prices rose 9.7\% in 2022,
that average conceals enormous variation. Fats and oils: +23.3\%.
Pasta and other grain products: +19.6\%. Fresh vegetables: +9.5\%.
Restaurant meals: +8.3\%.

As a consumer, you can partially escape measured food inflation through
\textit{substitution} --- exactly the behaviour that creates the CPI's
substitution bias. Switching from branded cereal to store-brand, choosing
chicken over beef, buying a larger format item or catching a sale --- all
of these reduce your \textit{personal} inflation rate below the measured
CPI for food.

This is not just a survival tip. It illustrates a deep economic point:
the CPI measures the cost of a \textit{fixed} basket. Real households
are not fixed-basket consumers --- they respond to price signals. Your
ability to substitute gives you more real purchasing power than the
headline inflation rate implies. The gap between your personal inflation
rate and the CPI, averaged across all consumers, is precisely the
\textit{substitution bias} in CPI measurement discussed in
Section~\ref{sec:cpi_measurement}. Estimated at 0.3--0.5 percentage
points per year on average, this bias was likely larger in 2022 when
relative price changes were unusually large --- giving households more
scope for substitution than normal.
\end{everydaylife}

%==========================================================
\section*{Chapter Summary}
%==========================================================
\addcontentsline{toc}{section}{Chapter Summary}

\begin{itemize}
  \item \textbf{Inflation} is a sustained increase in the general price
        level; \textbf{deflation} is a sustained decrease. Low positive
        inflation (2\% in Canada) is preferred to zero because it provides
        a policy buffer, accommodates relative wage adjustments, and
        partially offsets CPI measurement biases.

  \item The \textbf{CPI} is a Laspeyres index measuring the cost of a
        fixed basket of goods and services. It contains upward biases
        (substitution, quality, new goods, outlet) of approximately
        0.3--0.5 percentage points per year.

  \item \textbf{Core inflation measures} (CPI-trim, CPI-median, CPI-common)
        remove volatile components to reveal the underlying trend. All three
        rose above 5\% in 2022, confirming that the inflation surge was
        broad-based, not merely energy-driven.

  \item The \textbf{Fisher equation} ($r \approx i - \pi$) distinguishes
        nominal from real interest rates. In 2021, the real overnight
        rate in Canada was approximately $-3.2\%$ --- the most stimulative
        monetary stance in modern Canadian history.

  \item The \textbf{AD-AS model} explains two types of inflation:
        \textit{demand-pull} (rightward AD shift, rising prices and output)
        and \textit{cost-push} (leftward SRAS shift, stagflation). Canada's
        2021--2023 surge involved both simultaneously.

  \item The \textbf{short-run Phillips Curve} ($\pi = \pi^e - \alpha(u-u_n)$)
        shows a negative relationship between inflation and unemployment.
        The \textbf{long-run Phillips Curve} is vertical at the natural
        rate $u_n$: monetary policy cannot permanently reduce unemployment
        below $u_n$.

  \item Canada's \textbf{inflation-targeting framework} (since 1991) targets
        2\% CPI inflation within a 1--3\% control range. The Bank of Canada
        conducts policy primarily through the overnight rate.

  \item Monetary policy transmits through four channels: credit cost, asset
        prices, exchange rate, and expectations. Full effects take
        18--24 months.

  \item The 2022--2024 \textbf{rate-hike cycle} raised the overnight rate
        475 basis points to 5.00\% in 16 months --- the fastest tightening
        in modern Canadian history. Inflation returned to target by 2024
        without a technical recession, but at a significant cost to
        mortgage holders.

  \item \textbf{Inflation is regressive}: low-income households experienced
        approximately 8.5\% effective inflation in 2022 vs.\ 6.8\% for
        the top quintile, because they spend disproportionately more on
        food, shelter, and energy.
\end{itemize}

%==========================================================
\section*{Key Terms}
%==========================================================
\addcontentsline{toc}{section}{Key Terms}

\begin{description}[style=nextline, leftmargin=3.5cm]
  \item[\term{CPI (Consumer Price Index)}] A Laspeyres price index measuring
        the cost of a fixed basket of goods; Canada's primary inflation measure.
  \item[\term{CPI-common}] Bank of Canada core measure using factor analysis
        to extract economy-wide price movements.
  \item[\term{CPI-median}] The price change at the 50th percentile of the
        weighted component distribution.
  \item[\term{CPI-trim}] Core inflation measure that removes the top and
        bottom 20\% of price-change outliers.
  \item[\term{Core inflation}] Inflation measures that exclude the most
        volatile CPI components to reveal the underlying trend.
  \item[\term{Cost-push inflation}] Inflation arising from a leftward shift
        of the SRAS curve (rising input costs, supply shocks).
  \item[\term{Deflation}] A sustained fall in the general price level.
  \item[\term{Demand-pull inflation}] Inflation arising from a rightward
        shift of the AD curve (excess demand).
  \item[\term{Ex ante real rate}] The real interest rate based on expected
        future inflation; $r^e = i - \pi^e$.
  \item[\term{Ex post real rate}] The real interest rate based on actual
        realised inflation; $r = i - \pi$.
  \item[\term{Fisher equation}] The relationship $r \approx i - \pi$
        linking real rates, nominal rates, and inflation.
  \item[\term{Inflation targeting}] A monetary policy framework committing
        to keep inflation near a numerical target.
  \item[\term{Laspeyres index}] A price index using base-period quantity
        weights; the methodology underlying Canada's CPI.
  \item[\term{Long-run Phillips Curve}] The vertical relationship between
        inflation and unemployment at the natural rate $u_n$.
  \item[\term{NAIRU}] Non-Accelerating Inflation Rate of Unemployment;
        the unemployment rate consistent with stable inflation.
  \item[\term{Nominal interest rate}] The stated rate on a financial
        instrument; not adjusted for inflation.
  \item[\term{Overnight rate}] The Bank of Canada's policy rate; the rate
        at which major financial institutions lend to each other overnight.
  \item[\term{Phillips Curve}] The relationship between inflation (or wage
        inflation) and unemployment.
  \item[\term{Real interest rate}] The nominal interest rate adjusted for
        inflation; $r = i - \pi$.
  \item[\term{Real wage}] Nominal wage deflated by the price level;
        measures purchasing power of labour income.
  \item[\term{Sacrifice ratio}] Cumulative excess unemployment required
        per percentage-point reduction in inflation.
  \item[\term{Short-run Phillips Curve}] The negative relationship between
        inflation and unemployment at a given inflation expectation.
  \item[\term{Stagflation}] Simultaneous high inflation and high unemployment,
        typically arising from an adverse supply shock.
  \item[\term{Substitution bias}] CPI overstatement due to ignoring
        consumer substitution away from goods with rising relative prices.
  \item[\term{Zero lower bound}] The practical floor on nominal interest
        rates (approximately 0\%); limits the downward scope of monetary
        stimulus.
\end{description}

%==========================================================
\section*{Review Questions}
%==========================================================
\addcontentsline{toc}{section}{Review Questions}

\begin{enumerate}

\item \textbf{(Concept)} Define inflation and explain why a 2\% inflation
target is preferred to 0\% by most central banks. In your answer, discuss
the zero lower bound and the measurement bias arguments.

\item \textbf{(Measurement)} Canada's CPI basket assigns a 30\% weight
to Shelter and 20\% to Transportation. In 2022, Shelter rose 6.9\% and
Transportation rose 10.6\% year-over-year.
\begin{enumerate}[(a)]
  \item What is the combined contribution of these two categories to
        headline inflation, assuming the remaining categories average
        6.5\%?
  \item If all other categories fell to zero growth, what would headline
        inflation be?
  \item What does this illustrate about the sensitivity of CPI to
        its largest components?
\end{enumerate}

\item \textbf{(Fisher Equation)} In January 2021, the Bank of Canada
overnight rate was 0.25\% and CPI inflation was 1.0\%. By December 2022,
the overnight rate had risen to 4.25\% and inflation was 6.3\%.
\begin{enumerate}[(a)]
  \item Calculate the real interest rate in January 2021 and December 2022.
  \item In which period was monetary policy more accommodative (stimulative)?
        Explain using the Fisher equation.
  \item How does your answer change the interpretation of the rate hike
        cycle as ``aggressive''?
\end{enumerate}

\item \textbf{(AD-AS)} Use the AD-AS model to explain Canada's 2022
inflation surge. Draw a diagram showing:
\begin{enumerate}[(a)]
  \item The starting equilibrium at potential output $Y_f$.
  \item The rightward shift of AD resulting from COVID fiscal stimulus.
  \item The leftward shift of SRAS resulting from supply chain disruptions.
  \item The new equilibrium showing both higher prices and a different
        level of output.
\end{enumerate}
Label all curves and equilibria clearly.

\item \textbf{(Phillips Curve)} Using Equation~\ref{eq:phillips_curve}
and estimates $\alpha = 0.5$ and $u_n = 5.8\%$:
\begin{enumerate}[(a)]
  \item What does the model predict for inflation if the unemployment
        rate is 4.8\% and expected inflation is 2\%?
  \item What does the model predict if unemployment is 7.3\% and
        expected inflation is 3\%?
  \item If the Bank of Canada wants to reduce inflation from 5\% to 2\%
        and the sacrifice ratio is 2.0, how much cumulative excess
        unemployment (above $u_n$) would be required?
\end{enumerate}

\item \textbf{(Policy)} Why does the Bank of Canada monitor core inflation
measures (CPI-trim, CPI-median) rather than simply responding to headline
CPI? Under what circumstances would a supply shock cause the Bank to
\textit{tolerate} higher headline inflation without raising rates?

\end{enumerate}

%==========================================================
\section*{Quantitative Problems}
%==========================================================
\addcontentsline{toc}{section}{Quantitative Problems}

\begin{enumerate}

\item \textbf{CPI Calculation.} A household's monthly expenditure basket
consists of four categories:

\begin{center}
\begin{tabular}{lrrr}
\toprule
Category & Weight (\%) & Price (Base Year) & Price (Current Year) \\
\midrule
Food     & 20 & \$600 & \$672 \\
Housing  & 35 & \$1{,}200 & \$1{,}272 \\
Transport& 25 & \$400 & \$460 \\
Other    & 20 & \$300 & \$318 \\
\bottomrule
\end{tabular}
\end{center}

\begin{enumerate}[(a)]
  \item Calculate the Laspeyres CPI for the current year
        (base year = 100) using Equation~\ref{eq:cpi_formula}.
        Show all steps.
  \item Calculate the inflation rate $\pi$ from base year to current year.
  \item Which category contributed most to the measured inflation rate?
        Calculate each category's contribution to $\pi$.
  \item Suppose the household switches entirely from Transport to a
        category priced at base-year levels. What is the household's
        \textit{true} cost-of-living increase? How does it compare
        to the Laspeyres CPI answer? What bias does this illustrate?
\end{enumerate}

\item \textbf{Fisher Equation Applications.}
\begin{enumerate}[(a)]
  \item A 5-year Government of Canada bond yields 4.8\% nominal.
        If expected inflation over the next 5 years is 2.2\%, what
        is the expected real yield? Is this above or below the
        historical average neutral real rate of approximately 1\%?
  \item A student takes a variable-rate student loan at prime + 1.5\%.
        In January 2020, the prime rate is 3.45\% (overnight + 2.2\%).
        In December 2022, the prime rate is 6.45\%. The student's income
        is \$42{,}000 per year. CPI in 2020 was 137.0; in 2022 it was 151.9.
        Calculate the real interest rate on the student loan in each year
        and the real value of \$42{,}000 income in each year (in 2020 dollars).
  \item Canada runs a current account deficit when nominal interest rates
        are below trading partners' rates (capital flows out). In 2020, the
        Bank of Canada cut rates to 0.25\% while the ECB kept rates at
        $-0.50\%$ and the Fed cut to 0--0.25\%. Who had the relatively
        highest real rates given the respective inflation differentials?
        What does this imply for capital flows?
\end{enumerate}

\item \textbf{Real Wages and Cost of Living.}
Average weekly earnings in Canada were \$1{,}020 in January 2020
and \$1{,}147 in December 2022. CPI (all-items, 2002=100) was
137.9 in January 2020 and 156.0 in December 2022.
\begin{enumerate}[(a)]
  \item Calculate real average weekly earnings in each period
        (in January 2020 dollars) using Equation~\ref{eq:real_wage}.
  \item Calculate the nominal and real growth rates of weekly earnings
        over the period.
  \item If low-income workers (bottom quintile) earned \$720/week in
        January 2020 and \$780/week in December 2022, but experienced
        an effective inflation rate of 13.8\% over that period, what
        was the change in their real wages?
  \item What economic term describes a situation where nominal wages
        rise but real wages fall?
\end{enumerate}

\item \textbf{Phillips Curve and Sacrifice Ratio.}
Suppose the short-run Phillips Curve for Canada is:
$\pi_t = \pi^e_t - 0.6(u_t - 5.8\%)$
\begin{enumerate}[(a)]
  \item If expected inflation is 2\% and actual inflation is 4.5\%,
        what does the Phillips Curve imply about the unemployment rate
        at that time?
  \item The Bank of Canada wishes to reduce inflation from 4.5\% to 2\%
        over three years. In each year, the Bank accepts unemployment
        1.2 percentage points above the natural rate. What is the implied
        sacrifice ratio? Is this within the range of estimates in the
        literature (1.5--3.0)?
  \item If inflation expectations become fully rational (people correctly
        anticipate the Bank's actions), what happens to the sacrifice ratio?
        What does this tell us about the value of central bank credibility?
\end{enumerate}

\item \textbf{Distributional Analysis.}
Use the following simplified data to analyse distributional inflation:

\begin{center}
\begin{tabular}{lrrrr}
\toprule
Category & CPI Weight & 2022 Rise & Q1 Budget Share & Q5 Budget Share \\
\midrule
Food       & 16\% & 9.7\% & 22\% & 10\% \\
Shelter    & 30\% & 6.9\% & 38\% & 22\% \\
Energy     & 8\%  &28.5\% & 12\% &  4\% \\
Transport  & 12\% & 7.8\% &  8\% & 16\% \\
Services   & 20\% & 5.5\% & 12\% & 30\% \\
Other      & 14\% & 5.2\% &  8\% & 18\% \\
\bottomrule
\end{tabular}
\end{center}

\begin{enumerate}[(a)]
  \item Calculate the headline CPI inflation rate using official
        CPI weights and 2022 price changes.
  \item Calculate the effective inflation rate for Q1 (lowest-income)
        households using their budget shares as weights.
  \item Calculate the effective inflation rate for Q5 (highest-income)
        households.
  \item By how many percentage points did Q1 experience higher inflation
        than Q5? What drives this difference?
  \item Propose one policy instrument that could specifically compensate
        low-income households for this burden. What would be its fiscal cost
        if Canada has approximately 4 million households in Q1,
        each spending \$40{,}000 per year?
\end{enumerate}

\item \textbf{Data Exercise.} Visit the Bank of Canada website
(\url{www.bankofcanada.ca}) and find:
\begin{enumerate}[(a)]
  \item The current overnight rate target.
  \item The most recent annual CPI inflation reading.
  \item Using these two numbers, calculate the current real overnight rate.
        Is monetary policy currently stimulative, neutral, or restrictive
        (assume the neutral real rate is approximately 1\%)?
  \item Download the CPI-trim and CPI-median series from the Bank of Canada's
        core inflation page. Report whether core measures are above or below
        the 2\% target. What does this imply about the appropriate monetary
        policy stance?
\end{enumerate}

\end{enumerate}

%==========================================================
\section*{Essay Questions}
%==========================================================
\addcontentsline{toc}{section}{Essay Questions}

\begin{enumerate}

\item \textbf{(600--800 words)} ``The Bank of Canada's aggressive rate-hike
cycle in 2022--2023 was necessary and ultimately successful.'' Using the
AD-AS model, the Fisher equation, and the Phillips Curve, evaluate this
claim. Your essay should discuss: (a) the causes of the 2022 inflation surge;
(b) the mechanism through which rate hikes reduce inflation; (c) the costs
of the tightening cycle; and (d) your overall assessment.

\item \textbf{(500--700 words)} The CPI reports a single number for inflation,
but inflation is experienced differently across income groups, age cohorts,
homeowners vs.\ renters, and variable-rate vs.\ fixed-rate mortgage holders.
Using data from this chapter, write an essay arguing that ``the CPI is a
measure of average inflation, but Canada needs complementary distributional
inflation measures to make good policy.'' What additional measures would you
recommend Statistics Canada publish?

\item \textbf{(400--600 words)} Japan maintained near-zero inflation (and
frequently deflation) for much of the 1990s--2010s despite aggressive monetary
stimulus. Canada experienced 8.1\% peak inflation in 2022 despite similar
initial monetary conditions (near-zero rates). Using the AS-AD framework,
propose at least two hypotheses that explain why the same ultra-loose monetary
policy produced such different inflation outcomes in the two countries. What
does this imply about the limits of the quantity theory of money?

\end{enumerate}

%==========================================================
\section*{Discussion Questions}
%==========================================================
\addcontentsline{toc}{section}{Discussion Questions}

\begin{enumerate}

\item Fatima, from our opening vignette, is a variable-rate mortgage holder
who experienced both high inflation (eroding her real income) and sharply
higher mortgage payments (from the rate hikes). Is she a net winner or loser
from the sequence of events? What if she had locked into a 5-year fixed rate
at 2\% in 2020 instead? What if she were a renter with no mortgage?

\item The Bank of Canada held rates at emergency levels (0.25\%) throughout
2021, even as inflation rose above 4\%, citing concerns about pandemic
economic scarring. Was this the right call? What information would you need
to evaluate the Bank's decision in real time, rather than in hindsight?

\item ``Inflation is a hidden tax on savings, a transfer from creditors to
debtors, and a regressive burden on low-income households.'' Do you agree?
Who are the winners and who are the losers from an unexpected inflation
surge? Does your answer change if the inflation is anticipated vs.\ unexpected?

\item Canada's CPI Shelter component includes mortgage interest costs, which
rose sharply as the Bank of Canada raised rates. Should mortgage interest
costs be in the CPI? What are the arguments for and against the current
Canadian methodology? How would switching to ``owners' equivalent rent''
(the US approach) have changed the reported inflation numbers in 2023?

\end{enumerate}
